\section{Rangkuman}
\label{sec:rangkuman-bab1}

Pada bab ini telah dibahas:

\begin{itemize}
    \item Buku ini bertujuan membuat mahasiswa mampu \emph{bekerja}
          pada level instruksi mesin: menjelaskan cara kerja
          komputer, membaca hubungan instruksi dengan perangkat
          keras, menulis program assembly, dan memperbaiki program
          level rendah.

    \item Empat kemampuan akhir buku dipetakan kepada empat CPMK,
          yang bersama-sama mendukung \CPL{02} dan \CPL{03}.

    \item Pembahasan sengaja dibatasi pada prosesor 8086 mode real
          16-bit. Prosesor ini dipilih karena paling bersih untuk
          mempelajari konsep dasar, dan konsep tersebut berlaku
          langsung pada keluarga x86 modern.

    \item OBE merancang pembelajaran secara mundur: dari CPL menuju
          CPMK, Sub-CPMK, asesmen, baru kemudian materi. Materi
          hanya dianggap perlu bila mendukung indikator yang dapat
          diukur.

    \item Terdapat \dec{14} Sub-CPMK yang dipetakan ke
          \dec{16} minggu perkuliahan. Setiap bab materi buku ini
          bersesuaian dengan satu atau dua minggu perkuliahan.

    \item Setiap bab materi memiliki kerangka baku: Sub-CPMK, tujuan
          pembelajaran, materi pokok, contoh program, praktikum,
          latihan, rangkuman, dan evaluasi.

    \item Penilaian mata kuliah berbobot \dec{70}\% praktik dan
          \dec{30}\% teori, sehingga portofolio program dan jurnal
          praktikum menjadi bukti utama pencapaian kompetensi.

    \item Tiga perangkat lunak yang dipakai sepanjang semester adalah
          Emu8086, GUI Turbo Assembler, dan DOSBox.
\end{itemize}

\begin{checklistbox}
Tandai setiap butir yang sudah Anda kuasai.

\begin{itemize}
    \item[\checkbox] Saya dapat menjelaskan tujuan buku ini dan empat
          kemampuan akhir yang hendak dicapai.
    \item[\checkbox] Saya dapat menjelaskan jalur CPL $\rightarrow$
          CPMK $\rightarrow$ Sub-CPMK $\rightarrow$ materi
          $\rightarrow$ asesmen.
    \item[\checkbox] Saya dapat menemukan bab yang bersesuaian dengan
          minggu perkuliahan tertentu.
    \item[\checkbox] Saya sudah memiliki struktur direktori kerja
          untuk menyimpan portofolio program.
    \item[\checkbox] Saya sudah memiliki jurnal praktikum dengan
          format yang benar.
    \item[\checkbox] Saya sudah menyusun target belajar pribadi yang
          terukur.
    \item[\checkbox] Saya mengetahui letak glosarium dan rubrik
          penilaian pada lampiran.
\end{itemize}

\vspace{0.4em}
Seluruh butir di atas bersifat prasyarat. Pastikan semuanya bertanda
sebelum memasuki Bab~\ref{chap:landasan-teori}.
\end{checklistbox}

\begin{infobox}
\textbf{Rujukan bab.} Kerangka OBE pada bab ini mengikuti struktur
silabus dan RPS mata kuliah. Untuk memperluas wawasan mengenai
organisasi komputer yang menjadi latar mata kuliah ini, lihat
\cite{stallings2016} dan \cite{tanenbaum2013}. Untuk gambaran
praktis mengenai cara kerja prosesor 8086, lihat \cite{intel1979}.
\end{infobox}
